%%%%%%%%%%%%%%%%%%%%%%%%%%%%%%%%%%%%%%%%%%%%%%%%%%%%%%%%%%%%%%%%%%%
%                                                                 %
%                            ABSTRACT                             %
%                                                                 %
%%%%%%%%%%%%%%%%%%%%%%%%%%%%%%%%%%%%%%%%%%%%%%%%%%%%%%%%%%%%%%%%%%%


%\renewcommand{\simplechapterdelim}{.}
%\simplechapter

\chapter*{Περίληψη}
\addcontentsline{toc}{chapter}{Περίληψη}
Η Βαθιά Ενισχυτική Μάθηση (Deep Reinforcement Learning, DRL) αποτελεί μία από τις πλέον υποσχόμενες προσεγγίσεις για την ανάπτυξη αυτόνομων 
συστημάτων λήψης αποφάσεων σε σύνθετα και αβέβαια περιβάλλοντα. Στον τομέα των χρηματοοικονομικών συναλλαγών, οι πράκτορες DRL έχουν τη 
δυνατότητα να αλληλεπιδρούν δυναμικά με την αγορά και να βελτιστοποιούν άμεσα τη διαδικασία διαχείρισης χαρτοφυλακίου μέσω της μεγιστοποίησης 
μιας συνάρτησης ανταμοιβής. Παρά τη σημαντική πρόοδο που έχει σημειωθεί τα τελευταία χρόνια, οι χρηματοοικονομικές αγορές εξακολουθούν να παρουσιάζουν 
σημαντικές προκλήσεις λόγω της υψηλής μεταβλητότητας, της μη στασιμότητας, της μερικής παρατηρησιμότητας και των σημαντικών κοστών συναλλαγών. 
Επιπλέον, οι περισσότερες υπάρχουσες προσεγγίσεις βασίζονται αποκλειστικά σε αριθμητικά δεδομένα αγοράς, παραβλέποντας πληροφορίες που προέρχονται 
από ειδησεογραφικές πηγές και το επενδυτικό κλίμα.

Η παρούσα διπλωματική εργασία διερευνά την ανάπτυξη ενός προσαρμοστικού πράκτορα χρηματοοικονομικών συναλλαγών που συνδυάζει τεχνικές Βαθιάς 
Ενισχυτικής Μάθησης με πληροφορία συναισθήματος προερχόμενη από χρηματοοικονομικά κείμενα. Για τον σκοπό αυτό υλοποιείται ένα ολοκληρωμένο 
πλαίσιο συναλλαγών βασισμένο στον αλγόριθμο Proximal Policy Optimization (PPO) και σε αναδρομική αρχιτεκτονική Long Short-Term Memory (LSTM), 
η οποία επιτρέπει την αξιοποίηση χρονικών εξαρτήσεων και την αντιμετώπιση της μερικής παρατηρησιμότητας των χρηματοοικονομικών αγορών. 
Παράλληλα, αναπτύσσεται μηχανισμός εξαγωγής σήματος συναισθήματος μέσω προσαρμογής μεγάλου γλωσσικού μοντέλου (LLM) με χρήση της τεχνικής Low-Rank 
Adaptation (LoRA), αξιοποιώντας χρηματοοικονομικές ειδήσεις και κείμενα.

Κεντρική συνεισφορά της εργασίας αποτελεί η ενοποίηση πληροφοριών αγοράς, συναισθήματος και περιορισμών διαχείρισης ρευστότητας σε ένα κοινό πλαίσιο 
λήψης αποφάσεων. Συγκεκριμένα, διερευνάται η αξιοποίηση του σήματος συναισθήματος τόσο ως πρόσθετου χαρακτηριστικού στον χώρο καταστάσεων όσο και ως 
μηχανισμού άμεσης προσαρμογής των ενεργειών του πράκτορα. Επιπλέον, εισάγεται μηχανισμός ελέγχου ρευστότητας βασισμένος σε Περιορισμένες Μαρκοβιανές 
Διαδικασίες Απόφασης (Constrained Markov Decision Processes, CMDPs) με χρήση δυναμικά εκπαιδευόμενου πολλαπλασιαστή Lagrange, επιτρέποντας την ταυτόχρονη 
βελτιστοποίηση απόδοσης και διαχείρισης κινδύνου.

Η αξιολόγηση πραγματοποιείται σε περιβάλλον συναλλαγών κρυπτονομισμάτων μεγάλης κλίμακας και περιλαμβάνει εκτενή πειραματική ανάλυση διαφορετικών 
αρχιτεκτονικών, συναρτήσεων ανταμοιβής, συνόλων χαρακτηριστικών, μηχανισμών ελέγχου έκθεσης, περιορισμών ρευστότητας και στρατηγικών βελτιστοποίησης 
υπερπαραμέτρων μέσω Combinatorial Purged Cross Validation (CPCV). Τα αποτελέσματα δείχνουν ότι οι συμβατικές υλοποιήσεις πρακτόρων PPO παρουσιάζουν 
συχνά φαινόμενα υπερπροσαρμογής ή παθητικής κατάρρευσης έναντι της στρατηγικής αγοράς και διακράτησης. Αντίθετα, η χρήση αρχιτεκτονικής PPO-LSTM, 
η υιοθέτηση ανταμοιβής ενεργής απόδοσης, η επιβολή περιορισμών ρευστότητας και η ενσωμάτωση σήματος συναισθήματος οδηγούν σε σταθερότερη συμπεριφορά 
και βελτιωμένες εκτός δείγματος επιδόσεις. Ιδιαίτερα, η συνδυαστική αξιοποίηση των περιορισμών ρευστότητας και της πληροφορίας συναισθήματος 
παράγει τον βέλτιστο πράκτορα της εργασίας, ο οποίος επιτυγχάνει ανώτερη απόδοση και αποτελεσματικότερη διαχείριση κεφαλαίου σε σχέση με τις βασικές προσεγγίσεις αναφοράς.

Τα ευρήματα της εργασίας υποδεικνύουν ότι η ενσωμάτωση εναλλακτικών πηγών πληροφορίας μέσω Μεγάλων Γλωσσικών Μοντέλων, σε συνδυασμό με κατάλληλα 
σχεδιασμένους μηχανισμούς περιορισμών και διαχείρισης κινδύνου, μπορεί να βελτιώσει ουσιαστικά την αποτελεσματικότητα και τη γενικευσιμότητα 
πρακτόρων Βαθιάς Ενισχυτικής Μάθησης σε πραγματικά χρηματοοικονομικά περιβάλλοντα.




 
\chapter*{Abstract}
\addcontentsline{toc}{chapter}{Abstract}

Deep Reinforcement Learning (DRL) is one of the most promising approaches for developing autonomous decision-making systems in complex and uncertain environments. 
In the field of financial trading, DRL agents are capable of dynamically interacting with the market and directly optimizing the portfolio management process 
by maximizing a reward function. Despite the significant progress made in recent years, financial markets continue to present significant challenges due to high 
volatility, non-stationarity, partial observability, and significant transaction costs. Furthermore, most existing approaches rely exclusively on numerical market 
data, overlooking information derived from news sources and the investment climate.

This thesis investigates the development of an adaptive financial trading agent that combines deep reinforcement learning techniques with sentiment information 
derived from financial texts. To this end, a comprehensive trading framework is implemented based on the Proximal Policy Optimization (PPO) algorithm and a Long 
Short-Term Memory (LSTM) recurrent architecture, which allows for the exploitation of temporal dependencies and addresses the partial observability of financial 
markets. At the same time, a mechanism for extracting sentiment signals is being developed by fine-tuning a large language model (LLM) using the Low-Rank Adaptation 
(LoRA) technique, leveraging financial news and texts.

The key contribution of this study is the integration of market information, sentiment, and liquidity management constraints into a common decision-making framework. 
Specifically, the study explores the use of the sentiment signal both as an additional feature in the state space and as a mechanism for directly adjusting the 
agent’s actions. Furthermore, a fluidity control mechanism based on Constrained Markov Decision Processes (Constrained Markov Decision Processes, CMDPs) 
using a dynamically trained Lagrange multiplier, enabling the simultaneous optimization of performance and risk management.

The evaluation is conducted in a large-scale cryptocurrency trading environment and includes extensive experimental analysis of different architectures, 
reward functions, feature sets, exposure control mechanisms, liquidity constraints, and hyperparameter optimization strategies using Combinatorial Purged Cross 
Validation (CPCV). The results show that conventional PPO agent implementations often exhibit overfitting or passive collapse compared to the buy-and-hold strategy. 
In contrast, the use of a PPO-LSTM architecture, the adoption of an active return reward, the enforcement of liquidity constraints, and the incorporation of sentiment 
signals lead to more stable behavior and improved out-of-sample performance. In particular, the combined use of liquidity constraints and sentiment information 
produces the optimal trading agent, which achieves superior performance and more efficient capital management compared to the baseline approaches.

The findings of this study suggest that the integration of alternative information sources through Large Language Models, in combination with appropriately 
designed constraint and risk management mechanisms, can substantially improve the effectiveness and generalizability of Deep Reinforcement Learning agents 
in real-world financial environments.